\documentclass[letterpaper,11pt]{article}
\usepackage[utf8]{inputenc}
\usepackage[margin=1in]{geometry}
\usepackage{amsmath}
\usepackage{hyperref}

\title{Machine Learning-Enhanced Construction and Backtesting of Barra CN5 Style Factors in the A-Share Market
\\ {\small EN.553.640(01\&02) Machine Learning in Finance Final Project Report}} 
\author{
Sibo Tao
}  
\date{\today}

\begin{document}

\maketitle

\section{Introduction}
This final report extends the intermediate report and presents a complete machine-learning-enhanced quantitative research pipeline for A-share factor investing. The project objective is to bridge classical Barra-style factor investing and practical machine learning stock selection under realistic data engineering constraints.

Traditional multi-factor models are interpretable and robust in many institutional settings, but they usually assume linear mappings from factor exposures to returns and often ignore regime dependence, nonlinear interactions, and implementation frictions. In contrast, modern machine learning methods can model nonlinear structure and cross-factor interactions, yet they are often weak in economic interpretability and vulnerable to data leakage or overfitting if the pipeline is not designed carefully.

To address these issues, this project builds a modular architecture with strict anti-look-ahead rules, reusable intermediate artifacts, and unified evaluation outputs. The full system supports (i) raw data cleaning and panel construction, (ii) factor engineering and cross-sectional preprocessing, (iii) single-factor evaluation, (iv) machine-learning factor synthesis, and (v) backtesting and visualization. This design allows direct comparison between traditional hand-crafted factors and model-generated factors using the same IC/RankIC and quantile long-short framework.

The project also explicitly considers practical constraints: large-scale dataset size, long training time, and limited local compute resources. Therefore, this report not only summarizes what has been implemented, but also discusses optimization pathways and expansion plans for future work.

\section{Problem Setup}
The project studies stock selection in the Chinese A-share market. Let $i$ index stocks and $t$ index trading dates. For each stock-date pair, we construct standardized factor exposures
\begin{equation}
\mathbf{Z}_{i,t} = (Z_{i,t}^{(1)},\ldots,Z_{i,t}^{(K)})^\top,
\end{equation}
and define forward return labels
\begin{equation}
y_{i,t}=R_{i,t+h},
\end{equation}
where $h$ is the holding horizon (default: one day in the current implementation).

The core prediction/evaluation tasks are:
\begin{itemize}
    \item \textbf{Single-factor predictability}: evaluate each $Z^{(k)}$ via cross-sectional IC and RankIC.
    \item \textbf{Portfolio construction}: sort by factor/model score into quantiles and compute long-short return.
    \item \textbf{ML factor synthesis}: train model $g_\theta(\mathbf{Z}_{i,t})$ and treat output score $s^{ML}_{i,t}$ as a synthesized factor.
\end{itemize}

Major assumptions include: (1) announcement-date alignment or conservative lagging for financial statement fields, (2) tradable-universe filtering, (3) no usage of future information at rebalance date, and (4) forward labels with delisting-aware treatment whenever available.

\section{Theory}
\subsection{Theoretical Contribution}
Theoretical foundation comes from the multi-factor return decomposition
\begin{equation}
r_{i,t}=\sum_{k=1}^{K}X_{i,k,t}f_{k,t}+u_{i,t},
\end{equation}
and its cross-sectional interpretation: expected return depends on interpretable factor exposures. In this project, the key theoretical contribution is not a new theorem but an integrated formulation that links:
\begin{enumerate}
    \item classical factor pricing intuition,
    \item cross-sectional normalization/neutralization,
    \item and supervised ML mapping from factor vectors to future returns/ranks.
\end{enumerate}

This integrated view is significant because it preserves economic structure (factors, industry/size effects, implementable portfolios) while enabling nonlinear learning.

\subsection{Sketches of Proof}
Theoretical rationale can be summarized as follows. If factor preprocessing is cross-sectionally robust (winsorization, neutralization, standardization) and labels are forward-looking without leakage, then the estimated mapping from features to labels approximates an out-of-sample conditional ranking function rather than in-sample noise. Empirically, a useful approximation should produce:
\begin{itemize}
    \item positive mean IC or RankIC,
    \item stable quantile monotonicity,
    \item and economically interpretable long-short payoff after turnover constraints.
\end{itemize}

Although no formal consistency theorem is derived in this implementation report, the design follows standard empirical asset pricing and rolling-window validation logic under non-stationary financial time series.

\section{Practice}
\subsection{What the Project Implemented}
This project has implemented a full research and backtesting structure:
\begin{enumerate}
    \item \textbf{Data layer}: robust daily panel loading/cleaning, stock metadata integration, delisting metadata handling, and cached artifacts.
    \item \textbf{Universe and labels}: eligible universe generation and forward-return label construction.
    \item \textbf{Factor layer}: momentum, volatility, turnover, MACD-family and related price-volume factors with cross-sectional preprocessing.
    \item \textbf{Evaluation layer}: IC/RankIC time series, quantile backtest, long-short NAV and summary metrics.
    \item \textbf{ML extension}: rolling-window machine learning models (regression/classification/PCA-based synthesis) that output ML scores as factor files.
    \item \textbf{Visualization/reporting}: automated plotting and manifest-style output summaries.
\end{enumerate}

\subsection{Data}
The dataset consists of market daily bars, stock metadata, and currently integrated factor-related tables. It is already large enough to create substantial I/O and training pressure.

In addition to the intermediate-report data description, the reproducibility instructions are:
\begin{itemize}
    \item Data source (Baidu Netdisk): \url{https://pan.baidu.com/s/1Le6Z5yJciVsmWNzqb5wdoA?pwd=3iy8}
    \item Compressed size: about 3.11GB; extracted size: about 9.8GB.
    \item Rename the Chinese folder name to \texttt{A\_Share\_Data}.
    \item Decompress \texttt{daily\_hfq} inside the extracted directory.
    \item Place data under project root path: \texttt{data/raw/} (create both levels if missing).
\end{itemize}

\textbf{Project code repository: \url{https://github.com/NianXi1020/Quant_Finance_Factors_Backtesting_Structure/tree/codex/create-modular-data-loading-pipeline}}

\subsection{Experimental Setup and Results Summary}
The pipeline has been run in modular stages and produced:
\begin{itemize}
    \item shared stage artifacts (daily panel, metadata, universe, forward labels),
    \item ML feature panel with multi-million-row scale,
    \item rolling training outputs and factor parquet artifacts,
    \item evaluation and visualization outputs.
\end{itemize}

Main empirical observation so far: the framework is functionally complete for Stage-1 momentum/price-volume research and ML factor synthesis, but runtime is long under daily rebalancing and large sample horizons.

\section{Reflections}
\subsection{Theory}
Strengths of the current theoretical framing are interpretability and compatibility with established Barra-style research workflow. The major limitation is that no explicit finite-sample theorem is provided for ML-factor stability under non-stationarity, which remains an open challenge in practical quant research.

\subsection{Practice}
The main practical success is engineering completeness: the project can run from raw data to factor evaluation end-to-end. The key practical bottleneck is computational efficiency. Daily rolling retraining on large cross-sections is expensive, making experimentation cycles slow and hyperparameter tuning limited.

\section{Conclusion}
This final report presents a practical and extensible A-share factor research system that unifies classical factor modeling and machine-learning synthesis. The project demonstrates that robust data engineering, strict anti-leakage design, and unified evaluation are essential for meaningful ML-in-finance experiments. While Stage-1 implementation already supports strong baseline research, larger-scale factor expansion and compute-aware optimization are necessary for the next stage.

\section{Future Outlook Beyond Current Scope}
This section extends beyond the template structure and highlights next-step development priorities.

\subsection{Computation-Time Optimization for ML}
Current ML workflows are time-consuming due to large data volume and rolling retraining frequency. Future optimization priorities include:
\begin{itemize}
    \item reduce retraining frequency (e.g., monthly/weekly rebalancing for selected experiments),
    \item accelerate feature filtering and date indexing,
    \item cache intermediate model-ready matrices,
    \item parallelize model loops and tune resource allocation,
    \item profile hotspots and optimize high-cost data joins.
\end{itemize}

These optimizations are necessary to improve iteration speed, support broader hyperparameter search, and reduce engineering overhead.

\subsection{Data Expansion Under Better Compute Support}
The currently usable dataset is already around 9GB+ after extraction, and the present project timeline plus local compute limits constrained deeper expansion. With better computational resources, the next phase will:
\begin{itemize}
    \item integrate richer financial statement datasets,
    \item construct more traditional fundamental factors (value, quality, growth, leverage, accruals, efficiency),
    \item rebuild ML features on the expanded dataset,
    \item compare traditional-factor and ML-synthesized-factor performance under unified evaluation,
    \item incorporate turnover/cost-aware objectives into model selection.
\end{itemize}

In short, the current project established the pipeline backbone; the next phase will focus on richer factor breadth and stronger ML-depth comparisons under improved compute conditions.

\section*{Appendix}
The appendix content follows the same structure and reference basis as the intermediate report, including the same bibliography family and supporting mathematical notation conventions.

\begin{thebibliography}{99}

\bibitem{huatai_mfm_2016}
Huatai Securities Research Institute.
\textit{Huatai Multi-Factor Series I: A Preliminary Study of the Huatai Multi-Factor Model System}.
2016.

\bibitem{huatai_ai_2017_1}
Huatai Securities Research Institute.
\textit{Huatai Artificial Intelligence Series I: Overview of Artificial Intelligence Stock Selection Frameworks and Classical Algorithms}.
2017.

\bibitem{huatai_ai_2017_2}
Huatai Securities Research Institute.
\textit{Huatai Artificial Intelligence Series II: Generalized Linear Models for Artificial Intelligence Stock Selection}.
2017.

\bibitem{huatai_ai_2017_3}
Huatai Securities Research Institute.
\textit{Huatai Artificial Intelligence Series III: Support Vector Machine Models for Artificial Intelligence Stock Selection}.
2017.

\bibitem{huatai_ai_2017_4}
Huatai Securities Research Institute.
\textit{Huatai Artificial Intelligence Series IV: Naive Bayes Models for Artificial Intelligence Stock Selection}.
2017.

\bibitem{huatai_ai_2017_5}
Huatai Securities Research Institute.
\textit{Huatai Artificial Intelligence Series V: Random Forest Models for Artificial Intelligence Stock Selection}.
2017.

\bibitem{huatai_ai_2017_6}
Huatai Securities Research Institute.
\textit{Huatai Artificial Intelligence Series VI: Boosting Models for Artificial Intelligence Stock Selection}.
2017.

\bibitem{huatai_ai_2017_7}
Huatai Securities Research Institute.
\textit{Huatai Artificial Intelligence Series VII: Python Practice for Artificial Intelligence Stock Selection}.
2017.

\bibitem{huatai_ai_2017_8}
Huatai Securities Research Institute.
\textit{Huatai Artificial Intelligence Series VIII: Fully Connected Neural Networks for Artificial Intelligence Stock Selection}.
2017.

\bibitem{huatai_ai_2017_9}
Huatai Securities Research Institute.
\textit{Huatai Artificial Intelligence Series IX: Recurrent Neural Network Models for Artificial Intelligence Stock Selection}.
2017.

\bibitem{huatai_ai_2018_10}
Huatai Securities Research Institute.
\textit{Huatai Artificial Intelligence Series X: Applying Macro-Cycle Indicators to Random-Forest Stock Selection}.
2018.

\bibitem{huatai_ai_2018_12}
Huatai Securities Research Institute.
\textit{Huatai Artificial Intelligence Series XII: Feature Selection for Artificial Intelligence Stock Selection}.
2018.

\bibitem{huatai_ai_2018_13}
Huatai Securities Research Institute.
\textit{Huatai Artificial Intelligence Series XIII: Improving Loss Functions for Artificial Intelligence Stock Selection}.
2018.

\end{thebibliography}

\end{document}
